% Todd-corrected HKR proof of Poisson first-order vanishing.
\section{Todd補正から見た三次三様体の一次Poisson消滅}
\label{pf:sec:first-order}
中心三次対の曲率付きHochschild cochainによる直接的な一次MF比較とは別に、
本節では\emph{滑らかな三次三様体}について
Todd補正付きHKRとKuznetsov成分への制限を用いてPoisson方向の一次消滅を示す。
純粋なPoisson方向はKuznetsov成分の一次変形では消えるため、
可換模型が存在するなら、その方程式の本質的な変化は二次以降に現れる。
ここで証明するのは一次接写像の消滅であって、可換模型の存在、
残余smoothness (SM)、全次数での局所自明性ではない。

用いる外部結果は、Calaque--Rossi--Van den BerghのTodd補正付きHKRと
加群構造の両立性\cite[Theorem 1.4]{PFCRVdB}、および
Jacovskis--Lin--Liu--ZhangのHochschild作用の可換図式と
cubic threefoldに対する単射性\cite[Theorems 3.11 and 4.4(2)]{JLLZER}である。
以下の消滅命題はこれらの既知結果の適用であり、その帰結としての新規性は未確認である。

\subsection{通常のHKRとTodd補正の比較}
$X\subset\PP^4_{\CC}$を滑らかな三次超曲面とし、
\[
 \operatorname{HT}^2(X)
 =H^0(X,\wedge^2T_X)\oplus H^1(X,T_X)\oplus H^2(X,\OO_X)
\]
と書く。通常のHKR同型を
$I_{\mathrm{HKR}}:\operatorname{HT}^2(X)\xrightarrow{\sim}\HH^2(X)$、
加群構造と両立するTodd補正付き同型を$I^K$とする。
純粋なPoisson方向とは、一次変形類が
$I_{\mathrm{HKR}}(\pi)$、$\pi\in H^0(X,\wedge^2T_X)$、
$[\pi,\pi]=0$であることをいう。
同時に動かす可換複素構造の成分をこの条件に含めない。

\begin{lemma}[bivectorによる第一Chern類の縮約]
\label{pf:lem:chern}
滑らかな複素代数多様体$Z$と
$\pi\in H^0(Z,\wedge^2T_Z)$に対し、
\[
 \pi^\sharp c_1(K_Z)=0\quad\text{in }H^1(Z,T_Z).
\]
この主張にPoisson条件$[\pi,\pi]=0$は不要である。
\end{lemma}
\begin{proof}
$K_Z$を自明化するアフィン開被覆$\{U_i\}$と、
各$U_i$で消えない最高次微分形式$\mu_i$を取る。
重なりで$\mu_j=g_{ij}\mu_i$、$g_{ij}\in\OO^\times(U_i\cap U_j)$と書く。
$c_1(K_Z)$は、定数倍と符号の規約を除き、
$\{d\log g_{ij}\}$で表される$H^1(\Omega_Z^1)$の類である。

局所ベクトル場$D_i$を
\[
 \iota_{D_i}\mu_i=d(\iota_\pi\mu_i)
\]
で定める。最高次微分形式との縮約は
$T_{U_i}\simeq\Omega^{\dim Z-1}_{U_i}$という同型なので、$D_i$は正則である。
積の微分により
\begin{align*}
 \iota_{D_j-D_i}\mu_i
 &=g_{ij}^{-1}d\bigl(g_{ij}\iota_\pi\mu_i\bigr)
       -d(\iota_\pi\mu_i)\\
 &=d\log g_{ij}\wedge\iota_\pi\mu_i.
\end{align*}
右辺は、縮約と$\pi^\sharp$の符号を固定すれば
$\pm\iota_{\pi^\sharp(d\log g_{ij})}\mu_i$である。
従って$\{\pi^\sharp(d\log g_{ij})\}$は
$\{\pm D_i\}$の\v{C}ech coboundaryであり、そのコホモロジー類は零となる。
符号の選択は消滅の結論に影響しない。
\end{proof}

\begin{lemma}[cubic threefoldの純bivector上ではTodd補正が消える]
\label{pf:lem:todd}
任意の$\pi\in H^0(X,\wedge^2T_X)$について
\[
 I^K(\pi)=I_{\mathrm{HKR}}(\pi).
\]
\end{lemma}
\begin{proof}
\cite[Theorem 1.4]{PFCRVdB}の規約では
$I^K=I_{\mathrm{HKR}}\circ\iota_{\sqrt{\operatorname{td}(X)}}$である。
\[
 \sqrt{\operatorname{td}(X)}
 =1+\tfrac14c_1(T_X)+\theta_2+\cdots,
 \qquad\theta_2\in H^2(X,\Omega_X^2)
\]
を$\pi$に縮約すると、第一補正は
$\tfrac14\pi^\sharp c_1(T_X)\in H^1(T_X)$であり、
$c_1(T_X)=-c_1(K_X)$と補題\ref{pf:lem:chern}から零である。
第二補正は$H^2(X,\OO_X)=0$に値を取る。
三次以上の形式はbivectorには縮約できない。
よって補正は全て消える。逆Todd類を用いる規約でも同じ理由で結論は変わらない。
\end{proof}

\subsection{残余圏への制限写像の核}
$\mathcal K=\Ku(X)$とし、enhancementの完全忠実な包含から誘導される写像を
\[
 \operatorname{res}_X:\HH^2(X)\longrightarrow\HH^2(\mathcal K)
\]
と書く。

\begin{proposition}[Poisson方向の一次消滅]
\label{pf:prop:vanishing}
滑らかなcubic threefold $X$に対し、
\[
 \operatorname{res}_X\circ I_{\mathrm{HKR}}
 \big|_{H^0(X,\wedge^2T_X)}=0.
 \tag{PV}\label{pf:eq:vanishing}
\]
従って、実在する純Poisson一次変形からenhanced残余変形が誘導され、
そのKodaira--Spencer類が上の制限写像で与えられるなら、
残余変形の$\HH^2(\mathcal K)$における一次類は零である。
\end{proposition}
\begin{proof}
\cite[Theorem 3.11]{JLLZER}によれば、
$\alpha'=\operatorname{res}_X\circ I^K$と
\[
 \gamma:\HH^2(\mathcal K)\longrightarrow
 \operatorname{Hom}_{\CC}
 \bigl(H^1(X,\Omega_X^2),H^2(X,\Omega_X^1)\bigr)
\]
の合成はpolyvectorの縮約作用である。
ここで同論文のHochschild homologyの符号では、
$H\Omega_{-1}(X)=H^1(X,\Omega_X^2)$、
$H\Omega_1(X)=H^2(X,\Omega_X^1)$となる。
これら以外の該当するHodge成分は、cubic threefoldのLefschetz定理と
$H^i(X,\OO_X)=0$（$i>0$）により零である。

$\pi\in H^0(X,\wedge^2T_X)$の縮約は
\[
 H^1(X,\Omega_X^2)\xrightarrow{\ \iota_\pi\ }
 H^1(X,\OO_X)=0
\]
なので$\gamma\alpha'(\pi)=0$である。
cubic threefoldについて$\gamma$は単射である
\cite[Theorem 4.4(2)]{JLLZER}。従って$\alpha'(\pi)=0$を得る。
補題\ref{pf:lem:todd}で$I^K$を通常のHKRに戻せば
\eqref{pf:eq:vanishing}となる。
最後の主張は、仮定されたKodaira--Spencer類の制限との両立性による。
\end{proof}

\subsection{明示的復元式に対する制約}
\begin{corollary}[可換復元先の方程式は一次で動かない]
\label{pf:cor:second-order}
$X=V(F)\subset\PP^4$、$R=\CC[[h]]$とし、
特殊ファイバーが$\Ku(X)$であるenhanced残余族$\mathcal K_R$を考える。
その$R/(h^2)$への還元が、命題\ref{pf:prop:vanishing}の意味で
純Poisson一次変形から誘導されていると仮定する。
さらに滑らかな可換cubic族$Y_R=V(F_h)$とenhanced同値
\[
 \mathcal K_R\simeq\Ku(Y_R/R),\qquad Y_0=X
\]
が\emph{既に存在する}と仮定する。
$F_h=F+hG+O(h^2)$と書けば
\[
 [G]=0\quad\text{in }
 J(F)_3,
 \qquad J(F)=\CC[z_0,\ldots,z_4]/(\partial_0F,\ldots,\partial_4F).
 \tag{T0}\label{pf:eq:tangent-zero}
\]
従って$h=0$で恒等となる線形座標変換と方程式の単元倍により
$F_h=F+O(h^2)$とできる。
同じ結論は一次底$\CC[h]/(h^2)$上だけでも成立する。
\end{corollary}
\begin{proof}
命題\ref{pf:prop:vanishing}より、$\mathcal K_R$の一次類は零である。
enhanced同値はこの類を$\Ku(Y_R/R)$の一次類に送るので、
幾何的Kodaira--Spencer類$\kappa_Y\in H^1(X,T_X)$について
\[
 \eta(\kappa_Y)=0,\qquad
 \eta:H^1(X,T_X)\longrightarrow\HH^2(\Ku(X))
\]
となる。$H^1(T_X)$上のTodd補正も$H^2(\OO_X)=0$で消えるので、
ここでも通常のHKRと補正付きHKRは一致する。$\eta$は同型\cite[Theorem 4.4(2)]{JLLZER}であるから
$\kappa_Y=0$である。ここでは特殊ファイバーの同値が恒等である必要はなく、
Hochschild類の零は任意のenhanced同値で零に移る。

超曲面の法束完全列と射影空間のEuler列により、
$H^1(X,T_X)=J(F)_3$で、$\kappa_Y$は$[G]$に対応する。
従って一次形式$\ell_i$を用いて
$G=\sum_{i=0}^4\ell_i\partial_iF$と書ける。
座標変換$z_i\mapsto z_i-h\ell_i(z)$を代入すると一次項が消える。
方程式のスカラー変化もEuler恒等式
$\sum_i z_i\partial_iF=3F$によって同じJacobian部分に含まれる。
\end{proof}

\begin{remark}[一次消滅の適用条件と高次項]
\label{pf:rem:scope}
量子代数の交換子の半古典極限が$\pi$であることだけでは、
射影化した圏の一次変形が純粋に$I_{\mathrm{HKR}}(\pi)$であるとは断定できない。
可換方程式やgluingを一次で同時に動かすと、$H^1(T_X)$成分が加わり得る。
具体的な量子化族に系\ref{pf:cor:second-order}を適用する際は、
この純粋性とenhanced残余変形への制限との両立性を確認する。

確認後は、可換模型の最初の本質的な係数を
$F_h=F+h^2G_2+h^3G_3+\cdots$の$[G_2]\in J(F)_3$から探せばよい。
本節は$G_2$を計算しておらず、その非零性も、全ての奇数次項の消滅も主張しない。
また$\HH^2$での消滅はunderlying残余圏の一次類についての主張であり、
指定rotationを含む全coherentデータの自明化を結論しない。
存在・SM・高次同値の証明は、それぞれ独立に残る。
\end{remark}